\documentclass[11pt]{article}
\usepackage[margin=1in]{geometry}
\usepackage{amsmath,amssymb,amsthm,mathtools,hyperref,enumitem}
\usepackage[T1]{fontenc}
\usepackage{lmodern}
\hypersetup{colorlinks=true,urlcolor=blue,linkcolor=blue}
\newtheorem{theorem}{Theorem}[section]
\newtheorem{lemma}[theorem]{Lemma}
\newtheorem{proposition}[theorem]{Proposition}
\newtheorem{corollary}[theorem]{Corollary}
\newtheorem{definition}[theorem]{Definition}
\newcommand{\N}{\mathbb N}\newcommand{\Z}{\mathbb Z}\newcommand{\Q}{\mathbb Q}
\newcommand{\Tr}{\operatorname{tr}}\newcommand{\Res}{\operatorname{Res}}
\title{Branch-free algebraic observations as a reusable iPRA interface\\\large Finite multiplication matrices, elimination, and explicit proof obligations}
\author{Research note for the ProveIt iPRA development}
\date{16 November 2026}
\begin{document}\maketitle
\begin{abstract}
We propose a precise, reusable interface between constructive finitary mathematics and intuitionistic primitive recursive arithmetic (iPRA). Its distinguishing design choice is to represent a finite family of algebraic outcomes by a \emph{monic polynomial of observations}, rather than by a chosen root, branch, eigenvector, or labeled point. A polynomial observable is produced by a finite multiplication matrix, hence entirely by primitive-recursive arithmetic on coefficient codes, even when outputs coalesce. We prove the determinant/resultant presentation, polynomial substitution and tower-composition theorems, Newton trace identities, and a guard-dependent specialization theorem. We exhibit an explicit obstruction to selecting Puiseux coefficients from arbitrary Cauchy names and show exactly what our weaker interface does and does not circumvent. The accompanying Lean file is a specification and dependency map, \emph{not} a completed object-language derivation. The accompanying Python examples are exact-arithmetic tests, not iPRA proofs.
\end{abstract}
\tableofcontents

\section{Scope, repository baseline and the exact gap}
The repository \url{https://github.com/VladimirReshetnikov/ProveIt/tree/main/Logic/PrimitiveRecursiveArithmetic} distinguishes its checked \texttt{ConciseIPRA/} and \texttt{Derivations/} libraries from unverified research packets in \texttt{docs/proofs/}. In particular \texttt{Derivations/Core/Search.lean} gives an \emph{actual} object-language proof of bounded minimization, and \texttt{Core/LinearAlgebra.lean} defines finite matrix algebra. The packet \texttt{algebra/puiseux\_ipra/evidence/scope\_obstruction.md} establishes a nonuniformity obstruction for extracting Puiseux branches from arbitrary fast Cauchy coefficient names. The packet \texttt{foundations/ipra\_embedding\_framework\_2026-11-15} already develops finitary trace, guarded choice and certificate interfaces; and \texttt{foundations/ipra\_presentations\_to\_witnesses} already supplies a rational finite-action enumeration construction. We do \emph{not} claim to supersede these packages, or to derive Siegel's or Faltings' theorem.

Our specific contribution is a \emph{branch-free finite-observation algebra} whose constructors are coefficientwise elementary and whose compositional laws can be proved by finite identities. This reframes an otherwise discontinuous selection task as a stable, root-unlabeled calculation. It should be seen as a proposed common intermediate representation for the Puiseux, spectral-observable, algebraic-geometry and rational-point packets, rather than as a proof of their remaining analytic or height-bound steps.

\section{Primitive-recursive data and logical discipline}
Work initially over \(R=\Z\) or a computably presented \(\Q\), with canonical sign/denominator codes. A polynomial is a finite coefficient vector, a matrix is an array encoded by its dimensions and entries, and every bound on a loop is an explicit natural-number term. Matrix multiplication, determinants (by a fixed finite permutation sum), traces, monic reduction and bounded substitution are primitive-recursive \emph{schemata} on such codes. One must avoid asserting a single universal evaluator for all primitive-recursive functions: each fixed syntax tree gives its own symbol/definition in iPRA.

To internalize a lemma, replace the host mathematical notation with a first-order equality of coded primitive-recursive terms (or a Boolean equality test), expose all numerical parameters as free variables, and prove the universal open quantifier-free formula by the defining equations and quantifier-free induction. A host Lean proof about \(\Z\)-matrices does not establish an iPRA theorem until a term \texttt{Prf PRAAx} is built. We give a staged plan in Section~\ref{sec:lean}.

\section{Monic finite algebras and observation polynomials}
\begin{definition}[Monic algebra and multiplication observable]
Let \(f(Y)=Y^d+c_{d-1}Y^{d-1}+\cdots+c_0\in R[Y]\) be monic, \(d\ge1\). Set \(A=R[Y]/(f)\). Every class has a unique normal form of degree below \(d\) by repeated monic reduction. Let \(M_y\in R^{d\times d}\) denote multiplication by the class \(y\) of \(Y\), in the basis \(1,y,\ldots,y^{d-1}\). For \(g\in R[Y]\), define the \emph{observation polynomial}
\[
\mathcal O_f(g;T)=\det(TI_d-g(M_y))\in R[T].
\]
This definition never selects or orders roots.
\end{definition}

\begin{lemma}[Finite normal form]
For monic \(f\), every \(p\in R[Y]\) has a computable remainder \(r\) of degree \(<d\) and a quotient \(q\) such that \(p=qf+r\). The remainder is unique over every commutative ring \(R\), without assuming \(R\) is a domain.
\end{lemma}
\begin{proof}
For existence process coefficients from the leading term downward. Whenever a term \(aY^n\), \(n\ge d\), is present, subtract \(aY^{n-d}f\); its leading term cancels because \(f\) is monic. The maximum degree strictly decreases. Bounded polynomial input length supplies an explicit iteration count. For uniqueness, suppose \(r-r'=(q'-q)f\) with both \(r,r'\) of degree less than \(d\). If \(q'-q\) is nonzero, its highest nonzero coefficient equals the highest coefficient of its product with \(f\), because \(f\) is monic; hence the latter has degree at least \(d\). Contradiction. The zero ring is harmless: the statement is then immediate by direct coefficient equality. The argument uses only finite equality, not excluded middle for arbitrary propositions; on coded \(\Z\) coefficients equality is decidable.
\end{proof}

\begin{theorem}[Observation and resultant]
For all monic \(f\) and all \(g\in R[Y]\), \(\mathcal O_f(g;T)\) is monic of degree \(d\). It equals the determinant of the finite Sylvester-style map for \(f(Y)\) and \(T-g(Y)\), i.e. the monic \(f\) resultant
\[
\mathcal O_f(g;T)=\Res_Y(f(Y),T-g(Y)).\tag{1}
\]
In particular \(\mathcal O_f(Y;T)=f(T)\).
\end{theorem}
\begin{proof}
The matrix \(M_y\) is the companion matrix: the first \(d-1\) columns shift basis elements, and the last column is \((-c_0,\ldots,-c_{d-1})^\mathsf{t}\). Expansion of \(\det(TI-M_y)\) along the shift staircase gives \(f(T)\). The leading coefficient of \(\det(TI-g(M_y))\) is one, since exactly the identity permutation contributes \(T^d\).

For (1), perform monic division on the free module \(R[Y]_{<2d+\deg g}\), reducing multiples of \(f\). The determinant of the block change-of-basis matrix is one because the reduction is triangular with monic pivot entries. On the quotient free module \(A\), multiplication by \(T-g(y)\) is represented by \(TI-g(M_y)\). The usual Sylvester determinant, after these unit-triangular row/column operations, has this block and an identity block. The two determinants coincide. To avoid needing a variable-size Sylvester argument in the initial Lean milestone, equation (1) can instead be adopted as the \emph{definition} of the monic resultant; the equivalence to a classical resultant is then a later finite-matrix lemma. No argument about actual algebraic roots is used.
\end{proof}

\begin{theorem}[Finite observation composition]
Let \(g,h\in R[Y]\). The observation of \(h\) applied to the \(g\)-output is obtained by direct matrix evaluation:
\[
\mathcal O_f(h\circ g;U)=\det(UI_d-h(g(M_y))).\tag{2}
\]
If \(g(M_y)\) has distinct eigenvalues in a splitting extension, (2) is the product of \(U-h(g(\alpha_i))\) over roots with multiplicity; this explanatory splitting-field picture is \emph{not} required for its constructive proof. More generally, whenever \(f\mid k(g(Y))\), the explicit substitution map \(R[Z]/(k)\to R[Y]/(f)\), \(Z\mapsto g(y)\), is a well-defined ring homomorphism and commutes with polynomial observations on the target module.
\end{theorem}
\begin{proof}
Every polynomial evaluation in the matrix \(M_y\) is a finite sum/product, and \((h\circ g)(M_y)=h(g(M_y))\) by induction on the degree of \(h\). Taking characteristic determinants proves (2). For the quotient map, if \(p-p'=kq\), then \(p(g)-p'(g)=k(g)q(g)\) is divisible by \(f\); thus evaluation is independent of the representative. Addition and multiplication are preserved by their polynomial identities. Note that the observation \emph{multiplicities} belong to the chosen free target module: observations of a source quotient and target quotient are not automatically equal when their ranks differ. This avoids an invalid multiplicity-erasure step.
\end{proof}

\section{Symmetric trace calculus and division-free invariants}
\begin{theorem}[Newton identities without denominators]
Let \(B\) be any \(d\times d\) matrix over a commutative ring. Write
\[\det(TI-B)=T^d-e_1T^{d-1}+e_2T^{d-2}-\cdots+(-1)^de_d,\quad p_k=\Tr(B^k).\]
Then for \(1\le k\le d\),
\[
 p_k-e_1p_{k-1}+e_2p_{k-2}-\cdots+(-1)^{k-1}e_{k-1}p_1+(-1)^k k e_k=0.\tag{3}
\]
Thus all power-sum observations are finite polynomial expressions in elementary symmetric coefficients, with no division operation; conversely determining \(e_k\) from traces by division by \(k\) is \emph{not} valid over arbitrary rings.
\end{theorem}
\begin{proof}
Introduce commuting indeterminates \(z_1,\ldots,z_d\). The polynomial identity
\(E(t)=\prod_i(1+z_it)=\sum_{j=0}^d e_jt^j\) satisfies
\(tE'(t)=E(t)\sum_i\frac{z_it}{1+z_it}\) as a formal identity after multiplying by the denominator product. Comparing coefficients of \(t^k\), using the finite alternating expansion of each \((1+z_it)^{-1}\) only through degree \(k\), yields (3). Every coefficient is an integer polynomial, so the identity is valid after specialization to any commutative ring. For a matrix, the corresponding characteristic-polynomial and trace identities follow by specializing the universal generic matrix over \(\Z[b_{ij}]\), where the splitting principle is an optional explanation, not a necessary computation. Alternatively, prove the generic identity by the adjugate identity for \(I-tB\), taking a trace of its derivative; this is entirely polynomial. The resulting identity may be instantiated with integer-coded matrices and internalized by \texttt{ipra\_ring} plus bounded finite-sum lemmas.
\end{proof}

\section{The Puiseux obstruction and a branch-free repair}
Consider the monic quadratic family
\[f_a(Y)=Y^2-(t^2+a),\qquad a\ge0,\]
with \(t\) a formal parameter and \(a\) supplied by a fast Cauchy name. The Puiseux report shows that if an algorithm returned a chosen Puiseux expansion after \(t\mapsto t^e\), its coefficient at index \(e\) would separate \(a=0\) from \(a>0\), deciding the halting problem for a fixed primitive-recursive family of names.

\begin{theorem}[Continuous unlabeled polynomial observation]
In \(\Q[a,t][Y]/(Y^2-t^2-a)\), the multiplication matrix for \(y\) is
\[M_y=\begin{pmatrix}0&t^2+a\\1&0\end{pmatrix},\qquad
\mathcal O_{f_a}(y;T)=T^2-(t^2+a).\]
For the squared output,
\[\mathcal O_{f_a}(y^2;U)=(U-(t^2+a))^2.\tag{4}\]
The coefficients of both observation polynomials are polynomial functions of \(a,t\), hence admit uniformly computable fast Cauchy names from the supplied names of \(a,t\), on any fixed bounded input domain with an explicit precision-loss bound. No exact zero test on \(a\) is required.
\end{theorem}
\begin{proof}
The quotient relation gives \(y^2=t^2+a\), so multiplication by \(y\) interchanges basis vectors with the displayed coefficients. Its determinant is computed directly. Since multiplication by \(y^2\) is \((t^2+a)I_2\), (4) follows. On inputs \(|a|,|t|\le H\), polynomial evaluation is Lipschitz with an explicit rational constant depending on \(H\) and the finitely many coefficients: for monomials use \(|x^n-z^n|\le n\max(|x|,|z|)^{n-1}|x-z|\), then sum. Therefore a requested output accuracy \(2^{-n}\) reduces to input accuracy \(2^{-(n+c)}\) for an explicitly computable \(c\). This is a continuity statement about \emph{unlabeled finite polynomial data}; it does not supply separately named roots.
\end{proof}

\begin{proposition}[Boundary of the repair]
The observation interface cannot be extended uniformly to a branch selector that returns the individual coefficient at \(t^e\) in the above family from arbitrary Cauchy names. Conversely, given a \emph{supplied} exact algebraic presentation and explicit separation witnesses (e.g. a certified apartness bound and isolating discs), bounded search can select a unique isolated output and propagate that witness to a branch computation. The separation certificate is essential extra input, not an output of a discontinuous universal test.
\end{proposition}
\begin{proof}
For the negative part, suppose such a selector existed. Use the fixed primitive-recursive family
\(a_m=\sum_{k\ge0}2^{-k-1}H(m,k)\), where \(H(m,k)\) decides whether machine \(m\) halts by time \(k\). Its partial sums form uniformly fast Cauchy names. If \(a_m=0\), the \(t^e\) coefficient of either square-root branch has absolute value one; if \(a_m>0\), that coefficient is zero. An approximation within \(1/8\) therefore decides halting. This is impossible. For the positive conditional claim, a finite disjoint list of isolating boxes with rational separation \(\delta>0\) admits decidable point-to-box matching at precision less than \(\delta/4\); uniqueness follows from disjointness, and propagation is a finite computation. Actual convergence and isolation algorithms are additional domain-specific proof obligations, not consequences of the abstract interface.
\end{proof}

\section{A quantitative specialization rule}
\begin{theorem}[Guarded specialization of finite observations]
Fix degree \(d\), a polynomial expression \(g\), an integer coefficient-height bound \(H\), and a demanded output coefficient precision \(2^{-n}\). For any bounded family of monic polynomials of degree \(d\), there is an explicit primitive-recursive precision transformer \(\mu(d,H,\deg g,n)\) such that approximating every input coefficient within \(2^{-\mu}\) approximates every coefficient of \(\mathcal O_f(g;T)\) within \(2^{-n}\). On fixed \(d,\deg g\), this map can be computed solely by add/multiply error propagation through a fixed determinant circuit. To decide a sign or a zero/nonzero branch, however, one additionally needs an explicit gap \(|z|>2^{-q}\) or exact decidable algebraic input.
\end{theorem}
\begin{proof}
Normalize to a circuit using companion-matrix construction, Horner evaluation of \(g\), and determinant expansion as a sum of \(d!\) signed products. Assign to every gate both a rational magnitude bound \(B\) and an error bound \(\varepsilon\). For addition use \((B_1+B_2,\varepsilon_1+\varepsilon_2)\). For multiplication use \((B_1B_2, B_1\varepsilon_2+B_2\varepsilon_1+\varepsilon_1\varepsilon_2)\), with bounds enlarged by one to account for approximate inputs. The resulting error is an explicitly computed monotone polynomial in the common input error. Search a finite, explicitly overestimated exponent until the final error is below \(2^{-n}\); alternatively propagate dyadic exponents backward. Since each gate has finitely many inputs and multiplication depth is bounded by the circuit-size term, a crude \(\mu=n+O(L\log_2(2+H))\) with computable \(L\) suffices. No abstract continuity principle is invoked. Deciding \(z=0\) from Cauchy data is not justified by the error bound, whereas a supplied positive separation bound makes a finite comparison sound.
\end{proof}

\section{Applications and honest remaining obligations}
\paragraph{Puiseux and algebraic geometry.} Retain the equation or its finite algebra and compute polynomial observables, discriminants and elimination polynomials without splitting into labeled branches. At an explicitly certified simple-root region, augment the data with an isolating-box guard and obtain a local branch. At discriminant zero, do not silently continue that branch. This resolves the \emph{representation obstacle} for symmetric information, not the goal of uniformly returning all individual Puiseux coefficients.
\paragraph{Spectral observables.} For a finite matrix \(B\), use \(\det(TI-B)\), traces \(\Tr(B^k)\), and polynomial functional calculus to describe eigenvalues without selecting ordered eigenvectors. Gap-dependent spectral projectors require an additional certified spectral gap and inverse bounds, not just characteristic coefficients.
\paragraph{Rational points and height.} A finite multiplication algebra may supply annihilator identities and coordinate height bounds; the observations do \emph{not} turn finiteness alone into a height bound. The Faltings packet's missing general height theorem remains missing; the companion matrix constructor is only a finite algebraic subroutine.
\paragraph{Existing checked theorem bridge.} \texttt{Derivations/Core/LinearAlgebra}, \texttt{Core/MultiSums}, \texttt{Core/Permutation}, and \texttt{Core/Search} are natural insertion points. An object-level polynomial observable must be described by a \emph{fixed} PR symbol family and its defining equations, not by unrestricted quantification over all computation programs. For variable dimensions, use the existing bounded sums, products and matrix encodings; never assume quantifier-free induction over an unbounded existential predicate.

\section{Lean 4 formalization plan}\label{sec:lean}
The \texttt{lean/ObservationBlueprint.lean} file spells out a host-level interface with minimal definitions and theorem targets. It is deliberately noncompiled scaffolding; no claim of kernel checking is made. The completion path is:
\begin{enumerate}[leftmargin=*]
\item Define finite lists \texttt{coeff : Fin (d+1) -> ZT} with a forced monic top coefficient and an array/matrix evaluator over the \texttt{ZT} integer-difference encoding. Use \texttt{Core/LinearAlgebra} and \texttt{Core/Definitions} to implement the companion matrix and polynomial evaluation.
\item Establish companion identities by ring normalization and finite-sum reindexing. In particular first prove \(M_y^d+\sum_{i<d}c_iM_y^i=0\), then \(\det(TI-M_y)=f(T)\), using a finite determinant permutation expansion; rely on existing Cayley--Hamilton only if it actually has the requisite generic variable-dimension statement in \texttt{Prf PRAAx}.
\item Introduce fixed-index PR terms for \texttt{obsCoeff}, \texttt{polyEvalMatrix} and \texttt{matDet}, each via a primitive-recursive definition. For each use of a \texttt{Prf PRAAx} theorem prove their defining equations in the object system. Quantify over coefficient codes, dimension, and degree only when one PR coding and a single quantifier-free invariant are available.
\item Derive the term-level substitution identity \(\mathrm{eval}(h\circ g,M)=\mathrm{eval}(h,\mathrm{eval}(g,M))\) by quantifier-free induction on polynomial degree. Convert its equality into an observation equality by congruence of the determinant circuit.
\item Prove Newton's identities with bounded symmetric sums, induction and the integer-difference ring tool. Avoid dividing by \(k\). Lift the result to any coded integer matrix.
\item For precision transfer implement \texttt{boundGate} and \texttt{errorGate} by a fixed circuit evaluator and prove each local bound with rational cross-multiplication (no real-number object language required). The only induction is on finite circuit length.
\item Finally define object formula \texttt{ObsEq : Fm PTm n} and construct a term \texttt{Prf PRAAx (ObsEq)}, then invoke \texttt{pra\_iff\_native}. Explicitly audit that the theorem type is this judgment and that no imported host axiom is used as an object axiom.
\end{enumerate}
\paragraph{Verification status.} The current package supplies proofs in conventional mathematical language, a computational exact-arithmetic demonstration, and noncompiled Lean targets. The key remaining work is constructing the actual \texttt{Prf PRAAx} derivations for finite determinant polynomial identities and error propagation. This is a substantive theorem-proving task. This report makes no claim to have completed it.

\section{Sources and provenance}
\begin{enumerate}[leftmargin=*]
\item ProveIt iPRA root README: \url{https://github.com/VladimirReshetnikov/ProveIt/blob/main/Logic/PrimitiveRecursiveArithmetic/README.md}.
\item Checked derivation overview: \url{https://github.com/VladimirReshetnikov/ProveIt/blob/main/Logic/PrimitiveRecursiveArithmetic/Derivations/README.md}.
\item Existing finite-presentation framework: \url{https://github.com/VladimirReshetnikov/ProveIt/tree/main/Logic/PrimitiveRecursiveArithmetic/docs/proofs/foundations/ipra_embedding_framework_2026-11-15}.
\item Finite algebra rational point interface: \url{https://github.com/VladimirReshetnikov/ProveIt/tree/main/Logic/PrimitiveRecursiveArithmetic/docs/proofs/foundations/ipra_presentations_to_witnesses}.
\item Puiseux scope obstruction: \url{https://github.com/VladimirReshetnikov/ProveIt/blob/main/Logic/PrimitiveRecursiveArithmetic/docs/proofs/algebra/puiseux_ipra/evidence/scope_obstruction.md}.
\item Faltings scope warning: \url{https://github.com/VladimirReshetnikov/ProveIt/tree/main/Logic/PrimitiveRecursiveArithmetic/docs/proofs/diophantine-number-theory/faltings_ipra}.
\end{enumerate}
\end{document}
